\documentclass{jsarticle}
\usepackage{amsmath}
\usepackage{amssymb}
\begin{document}
\title{大域的微分方程式についてのレポート}
\author{Masaaki Yamaguchi, with my son}
\date{}
\maketitle

  非対称性理論によるパリティの破れから、宇宙と異次元にはそれぞれ、
  重力と反重力が存在しているが、ヒッグス場の方程式に宇宙での数式を入力すると、
  ゼータ関数だけが出力される。対称性の仕組みから異次元が
  対として生成される。
  これに関連するCP対称性の破れは、反粒子と粒子がそれぞれ、
  弦理論から結合している宇宙と異次元から、反粒子が消えたのは、
  結合した結果、宇宙のゼータ関数が出力されたためでもある。
  ヒッグス場の式から、宇宙と異次元は、それぞれ別に存在している。
  $ \Omega $におけるp形式による外微分が、
  $ \mathcal{L}_x (d \Omega) = d(\mathcal{L} \Omega) $と同値により、
  次の方程式が成り立つ。
  $$ {d \over df}F = e^{x \log x} $$
  $$ {d \over df}F = F^{{(f)}^{'}} $$
  $$ = m(x), x^{{1 \over 2} + iy} = e^{x \log x} $$
  この式からゼータ関数がガロワ拡大に導かれて、ヒッグス場の式と
  同値になる。
  $$ f(x) + g(y) \ge 2 \sqrt{f(x) g(y)} $$ 
  $$ {d \over df}F \ge {d \over df}\int \int{1 \over (x \log x)^2}dx_m 
  + {d \over df}\int \int{1 \over (y \log y)^{1 \over 2}}dy_m $$
  $$ {d \over df}  \int \int{1 \over (x \log x)^2}dx_m = \log (x \log x) $$
  $$ {d \over df}\int \int{1 \over (y \log y)^{1 \over 2}}dy_m = 
  2(y \log y)^{1 \over 2} $$
  $$  {d \over df} F(x_m) = {\partial \over \partial f}F(x_m),
  \left({\int \int{1 \over (x \log x)^2}dx_m}\right)^{(f)^{'}} $$ 
  $$ {d \over df} F(y_m) = {\partial \over \partial f}F(y_m),
  \left({\int \int{1 \over (y \log y)^{1 \over 2}dy_m}}\right)^{(f)^{'}} $$
  $$ \bigoplus {\mathcal{H}\Psi \over \nabla \mathcal{L}} 
     = {1 \over {{d \over d f}F}}  $$
  $$ = {d \over d \Lambda}\lambda $$
  ヒッグス場の式を逆関数で求めると、世界面を切断した値を場として、
  張力として求められる。この値を単体分割した$ \tau $を多様体にすると、
  多様体の計量で超関数を正規部分群で組み合わせ多様体として、
  ウィークスケールによるエネルギーにこの逆関数は至る。
  $$ T^{\mu\nu} =  {G^{\mu\nu} + {1 \over 2}g_{ij}\Lambda} $$
  $$ f = \tau = \left(T^{\mu\nu}\right)^{-1} $$
  $$ T^{\mu\nu} = \int \tau(x,y) dx_m dy_m $$
  $$ \bigoplus \left({\mathcal{H}\Psi}^{\nabla} \right)^{\oplus L} = 
  i \hbar \psi $$
  $$ = \bigoplus {{\left({i \hbar}^{\nabla} \right)}^{\oplus L}} $$
  微分幾何の量子化は、ハイゼンベルク方程式の確率振幅における、
  ウィークスケールによるエネルギー最小単位の
  エントロピー不変量でもある。ヒッグス場のエネルギーの確率振幅でもある。
  ヒッグス場のエネルギーが加群分解して再接続されてもいる。
  大域的微分方程式とも同型である。
  $$ F = {1 \over 4}x^4 + C, f = x^3, e^{i \theta} = \cos \theta +
  i \sin \theta $$
  $$ F = - \cos \theta, f = \sin \theta $$
  $$ {d \over df}F = {({1 \over 4}x^4 + C)}^{(x^3)}, 
  {d \over df}F = {(-\cos \theta)}^{(\sin \theta)} $$
  $$ {({1 \over 4}x^4 + C)}^{(x^3)} = {({1 \over 4}x^4 + C)}^{(x^3)^{'}} $$
  $$ {1 \over 4}(x^3)^{x^3} x + {1 \over 4}x^4 (x^x)^3 
  = {1 \over 4}e^{f \log f} + {1 \over 4}x^4 (e^{x \log x})^3 $$
  $$  = {1 \over 4}({{x^3}^{x^3}})^{'} \cdot {x^{(x \cdot x^3)}}^{'}  $$
  $$ = {1 \over 4}{x^{x^3 \log x^3}}^{'} \cdot {x^{x^4 \log x}}^{'} $$
  $$ = {3x^2 \over 4}e^{3x^2 \log x}\cdot (3x^2 \log x + 1) $$
  $$ {d \over df}F = {(-\cos \theta)}^{(\sin \theta)} = 	
   (- \sin \theta)^{{'}^{(\sin \theta)}} $$
  $$ = (f)^{{'}^{(f)}} = {((f)^f)}^{'} $$
  $$ = e^{x \log x} $$
\end{document}

